\chapter{Identity, Reuse, and What a Cache Key Cannot Be}

The most attractive idea in any fast-indexer project is that a second
scan should be faster than the first. The natural mechanism is a
cache of parsed subtree summaries: walk the FS tree once, summarise
each node, and on the next run reuse the summary if the node has not
changed. The chapter is about why the obvious cache key for that
mechanism is dangerous, and what the right shape of a safe key looks
like.

\section{The Tempting Mistake}

The most obvious cache key is the OID: it identifies an object across
versions, it is stable across copy-on-write rewrites, and it is small.
A summary cache keyed on \texttt{(volume\_uuid, oid)} would deduplicate
unchanged subtrees across runs.

It would also return stale summaries, silently, the first time the
volume mutates a subtree the cache already knows.

\textbf{Observation.} The FS-tree root's logical OID stays stable
across mutations; everything else about the on-disk object changes.
EX-06 mutated the proof fixture, recorded the FS-tree root identity
fields before and after, and observed:

\begin{center}
\begin{tabular}{ll}
\toprule
field & across mutation \\
\midrule
logical OID                & stable \\
physical address (paddr)   & changed \\
object XID                 & changed \\
object checksum            & changed \\
block hash                 & changed \\
\bottomrule
\end{tabular}
\end{center}

A cache keyed on OID would return the pre-mutation summary for the
post-mutation tree. The parser would emit rows derived from the old
state and aggregate them into the new state's directory totals. The
user would see a treemap that does not match the disk.

The OID is a \emph{virtual address}. It names \emph{some} version of
the object; it does not name a particular content. A cache that wants
content identity must use a content identifier.

\section{The Shape Of A Safe Cache Key}

A summary cache that is safe under mutation must distinguish entries
that share an OID but differ in content. The minimum sufficient key
is a tuple whose members come from several layers:

\begin{description}
  \item[Source identity.] A volume UUID, a container UUID, and an
    image fingerprint (for detached sources). Distinguishes one
    APFS volume from another, and distinguishes the same volume's
    different image dumps.
  \item[Parser version.] A monotonic identifier for the body decoder
    and the namespace builder. A parser change invalidates every
    summary, because a summary's interpretation depends on the
    parser that produced it.
  \item[Summary schema version.] A monotonic identifier for the
    summary structure itself. A schema change is independent of a
    parser change.
  \item[OMAP context.] Whether the OID was resolved through the
    container OMAP or the volume OMAP. A collision between the two
    name spaces is rare but representable.
  \item[OID.] The virtual identifier.
  \item[Object XID.] The version of the object the summary was built
    for.
  \item[Physical address.] Useful as a cross-check; an OMAP that
    resolves \texttt{(oid, xid)} to a different paddr than the cache
    expects invalidates the entry.
  \item[Checksum or content hash.] The strongest individual signal.
    A block whose checksum has changed must not reuse a summary
    keyed on the old checksum.
  \item[Expected type and subtype.] Defensive. An OID that has been
    re-allocated to a different object class is a new object.
\end{description}

The key is therefore not a small integer. It is a record of nine
fields, eight of which the parser already needs to compute for any
read, and the ninth (parser version) is constant per build.

\section{What EX-07 Proved And What It Did Not}

The project's first probe of the cache idea (EX-07) restricted itself
to a narrow question: under exact node-identity matches on every
field of the key above, does subtree reuse ever produce a wrong
summary?

\textbf{Observation.} EX-07 found zero false reuse for exact
node-identity matches in the detached lab corpus.

That is the strongest positive result a subtree-reuse cache can ask
for at this stage. It means the mathematical foundation is sound:
if two parses see the same node identity, they will get the same
summary, and reusing one for the other is safe.

It is not yet a production cache theorem. The probe ran on one
corpus, on detached images, with mutations limited to a small set.
Three things still have to land before a cache is safe to ship:

\begin{enumerate}
  \item Re-run EX-07's mutation grid against the native Rust full
    parse (now possible, since the parser and the parallel
    fallback of chapter~11 are both production code) to confirm
    the finding holds when the parser is this project's own
    production code rather than a third-party reference.
  \item Build a cache storage design with crash-safe writes; a
    cache that can corrupt during a kill produces stale-and-wrong
    rather than absent-and-correct, which is the failure mode the
    fail-closed discipline is supposed to prevent.
  \item Specify an invalidation policy that names every condition
    under which a summary becomes stale: parser version mismatch,
    missing volume continuity, source identity mismatch,
    unsupported feature drift, checkpoint rollback, snapshot
    switch, unknown record-body type, or any new metric the cache
    did not record (the allocated-size column is the most recent
    example of a column a pre-existing cache would not have
    recorded).
\end{enumerate}

Until those gates pass, the safe default is full reparse, validated
against the namespace, logical-size, and allocated-size oracles. A
second scan of the same volume produces the same answer as the
first; the parallel walker has made it produce that answer faster,
but it does not yet produce it incrementally.

\section{The Identity-Tracking Discipline}

A small detail of EX-06 deserves promotion. The probe did not assume
which identity fields would change; it recorded all of them, every
time. That habit --- "record every plausibly-load-bearing identity
field at every relevant boundary, before deciding which ones matter"
--- is what made the cache discussion concrete instead of speculative.

The same habit appears in EX-12's OMAP lookup contract: the probe
records every resolved \texttt{(oid, paddr, xid)} sample and replays
header validation at each, rather than checking only the resolutions
that the parser happens to read. The cost of recording is small;
the cost of discovering, later, that a key field was never captured
is large.

The discipline generalises beyond caching. A reverse-engineering
project produces evidence whose value depends on what it recorded
at the time. The right rule is: when the question is "which fields
matter," the experiment records all of them.

\section{What This Chapter Forecloses}

The chapter forecloses two specific design moves that other indexers
sometimes make.

\paragraph{No bare-OID cache.} The parser will never key reuse on
OID alone, even with parser-version qualification. The mutation
behaviour above means a bare-OID cache is unsafe in principle, not
just unsafe under unusual conditions.

\paragraph{No partial summary reuse across mismatched scan XIDs.}
A summary computed at scan XID $a$ is not reused at scan XID $b$.
The whole point of pinning the scan XID is that object identity
includes the XID; relaxing that turns the cache into a stale-state
producer. The cache may compare keys across runs only when the keys
agree on every field, scan XID included.

The chapter does not foreclose the cache idea itself. Subtree reuse
is a real opportunity --- the EX-07 result is encouraging --- but
the cost of getting it wrong is high enough that the project will
land it only after the gates above are met. Until then, every scan
fully reparses, and the manual records why.
